% Лекц 11: семинар ба гэрийн даалгавар. Бодлогууд: Г.Н. Берман, "Сборник задач по курсу
% математического анализа" (22-р хэвлэл). [Б.N] нь номын бодлогын дугаар.

\ProblemSet{Семинар 11}
\begin{frame}[allowframebreaks]{Семинар 11}
\pgroup{Ролль ба Лагранжийн теорем.}
\begin{plist}
\prob{1116}{$y=x^3+4x^2-7x-10$ функцийн хувьд $[-1,2]$ хэрчим дээр Роллийн теоремыг шалга.}
\prob{1117}{$y=\ln\sin x$ функцийн хувьд $\bigl[\frac{\pi}{6},\frac{5\pi}{6}\bigr]$ хэрчим дээр Роллийн теоремыг шалга.}
\prob{1131}{$y=\ln x$ функцийн хувьд $[1,e]$ хэрчим дээр Лагранжийн теоремыг шалга.}
\prob{1132}{Лагранжийн томьёогоор $\frac{a-b}{a}\le\ln\frac{a}{b}\le\frac{a-b}{b}$ \ ($0<b\le a$) тэнцэтгэл бишийг батал.}
\end{plist}
\pgroup{Монотон байх завсрыг ол.}
\begin{plist}
\prob{1143}{$y=2x^3+3x^2-12x+1$ функц $(-2,1)$ завсарт буурахыг харуул.}
\prob{1150}{$y=x^3-3x^2-9x+14$}
\prob{1157}{$y=x^2e^{-x}$}
\end{plist}
\pgroup{Экстремумыг ол.}
\begin{plist}
\prob{1165}{$y=2x^3-3x^2$}
\prob{1166}{$y=2x^3-6x^2-18x+7$}
\prob{1175}{$y=x-\ln(1+x)$}
\end{plist}
\pgroup{Хэрчим дээрх хамгийн их ба хамгийн бага утгыг ол.}
\begin{plist}
\prob{1185}{$y=x^4-2x^2+5$, \ $[-2,2]$}
\prob{1186}{$y=x+2\sqrt x$, \ $[0,4]$}
\end{plist}
\pgroup{Тэнцэтгэл бишийг батал.}
\begin{plist}
\prob{1199}{$e^x>1+x$ \ ($x\ne0$)}
\prob{1200}{$x>\ln(1+x)$ \ ($x>0$)}
\prob{1212}{$18\times18$ см хэмжээтэй дөрвөлжин картоны булангуудаас ижил дөрвөлжин хайчилж аваад, ирмэгийг нь нугалж таггүй хайрцаг хийнэ. Хайрцагны эзлэхүүн хамгийн их байхын тулд хайчлах дөрвөлжны тал ямар байх вэ?}
\end{plist}
\pgroup{Лопиталийн дүрмээр хязгаарыг ол.}
\begin{plist}
\prob{1325}{$\displaystyle\lim_{x\to0}\frac{\ln\cos x}{x}$}
\prob{1330}{$\displaystyle\lim_{x\to0}\frac{x-\sin x}{x-\tg x}$}
\prob{1346}{$\displaystyle\lim_{x\to+\infty}x^ne^{-x}$}
\prob{1349}{$\displaystyle\lim_{x\to1}\Bigl[\frac{x}{x-1}-\frac{1}{\ln x}\Bigr]$}
\prob{1358}{$\displaystyle\lim_{x\to0}x^{\sin x}$}
\end{plist}
\end{frame}

\ProblemSet{Гэрийн даалгавар 11}
\begin{frame}[allowframebreaks]{Гэрийн даалгавар 11}
\pgroup{Ролль ба Лагранжийн теорем.}
\begin{plist}
\prob{1118}{$y=4^{\sin x}$ функцийн хувьд $[0,\pi]$ дээр Роллийн теоремыг шалга.}
\prob{1119}{$y=\sqrt[3]{x^2-3x+2}$ функцийн хувьд $[1,2]$ дээр Роллийн теоремыг шалга.}
\prob{1124}{$x^3-3x+c=0$ тэгшитгэл $(0,1)$ завсарт хоёр ялгаатай язгууртай байж чадахгүйг харуул.}
\prob{1125}{$f(x)=(x-1)(x-2)(x-3)(x-4)$ функцийн уламжлалыг олохгүйгээр $f'(x)=0$ тэгшитгэл хэдэн бодит язгууртай, тэдгээр нь ямар завсруудад оршихыг тогтоо.}
\prob{1128}{$y=x(1-\ln x)$ функцийн хувьд $[a,b]$ хэрчим дээр Лагранжийн томьёог бич.}
\end{plist}
\pgroup{Монотон байх завсрыг ол.}
\begin{plist}
\prob{1144}{$y=\sqrt{2x-x^2}$ функц $(0,1)$-д өсөж, $(1,2)$-т буурахыг харуул.}
\prob{1145}{$y=x^3+x$ функц хаа сайгүй өсөхийг харуул.}
\prob{1146}{$y=\arctg x-x$ функц хаа сайгүй буурахыг харуул.}
\prob{1151}{$y=x^4-2x^2-5$}
\prob{1154}{$y=\frac{1-x+x^2}{1+x+x^2}$}
\prob{1156}{$y=x-e^x$}
\prob{1158}{$y=\frac{1}{\ln x}$}
\end{plist}
\pgroup{Экстремумыг ол.}
\begin{plist}
\prob{1161}{$y=2\sin x+\cos2x$ \ ($0\le x\le2\pi$)}
\prob{1167}{$y=\frac{3x^2+4x+4}{x^2+x+1}$}
\prob{1173}{$y=\frac{1+3x}{\sqrt{4+5x^2}}$}
\prob{1176}{$y=x-\ln(1+x^2)$}
\end{plist}
\pgroup{Хамгийн их ба хамгийн бага утгыг ол.}
\begin{plist}
\prob{1189}{$y=\sqrt{100-x^2}$, \ $-6\le x\le8$}
\prob{1191}{$y=\frac{x-1}{x+1}$, \ $0\le x\le4$}
\end{plist}
\pgroup{Тэнцэтгэл бишийг батал.}
\begin{plist}
\prob{1201}{$\ln x>\frac{2(x-1)}{x+1}$ \ ($x>1$)}
\prob{1202}{$2x\arctg x\ge\ln(1+x^2)$}
\end{plist}
\pgroup{Хамгийн их, хамгийн бага утгын бодлого.}
\begin{plist}
\prob{1208}{8 тоог хоёр нэмэгдэхүүнд хувааж, тэдгээрийн кубуудын нийлбэр хамгийн бага байхаар хуваа.}
\prob{1209}{Ямар эерэг тоог өөрийнх нь урвуу тоотой нэмэхэд хамгийн бага нийлбэр гарах вэ?}
\prob{1215}{Таггүй сав цилиндр хэлбэртэй, эзлэхүүн нь $v$. Гадаргуу нь хамгийн бага байхын тулд суурийн радиус ба өндөр ямар байх вэ?}
\prob{1221}{$R$ радиустай бөмбөрцөгт багтсан хамгийн их эзлэхүүнтэй цилиндрийн өндрийг ол.}
\end{plist}
\pgroup{Хязгаарыг ол.}
\begin{plist}
\prob{1326}{$\displaystyle\lim_{x\to0}\frac{e^x-1}{\sin x}$}
\prob{1328}{$\displaystyle\lim_{x\to0}\frac{x-\arctg x}{x^3}$}
\prob{1338}{$\displaystyle\lim_{x\to0}\frac{e^x-e^{-x}-2x}{x-\sin x}$}
\prob{1352}{$\displaystyle\lim_{x\to0}\Bigl(\ctg x-\frac1x\Bigr)$}
\prob{1360}{$\displaystyle\lim_{x\to0}\Bigl(\frac1x\Bigr)^{\tg x}$}
\prob{1361}{$\displaystyle\lim_{x\to\infty}(e^x+x)^{1/x}$}
\end{plist}
\end{frame}
